
\subsection{Les trois états de paramètre}

On considère que les données d'un paramètre peuvent passer par trois états
différents.
La construction du nom du vecteur contenant ces données dépend de cet état.
Ces états sont :
\begin{description}
\item[•] {brut (ou quanta)} : \\
         les données brutes d'un paramètre sont une simple et exacte copie du
         contenu du fichier d'origine (issu généralement du système
         d'acquisition). Elles n'ont subi aucune transformation
         lors de la lecture puis de la ré-écriture en vecteur. \\
         Le terme de quanta est appliqué aux données dont l'unité n'est pas
         une grandeur physique mesurée directement. Les données sont par
         exemple en Ohm et il faut leur appliquer un calcul purement
         automatique, pour connaitre la grandeur géophysique correspondante.
         \textbf{\textit{Note : la notion de données de type quanta tend a
         disparaître avec la modernisation des systèmes d'acquisition.}}

\item[•] {calibré} : \\
         les données calibrées d'un paramètre sont les données brutes
         auxquelles a été appliqué un calcul purement automatique, dans le but
         d'obtenir la grandeur physique mesurée (une température,
         une pression, etc.)

\item[•] {traité} : \\
         les données traitées d'un paramètre sont des données de n'importe quel
         type, qui ont fait l'objet d'une expertise et d'une validation
         humaine ou automatique.
\end{description}


\subsection{Construction du nom des vecteurs}
Le cas des données de type quanta est simple : leur nom du vecteur est celui
qui est trouvé dans le fichier d'origine, écrit entièrement en majuscules. \\

Pour tous les autres types de donnée (brut, calibré, traité), le schéma
de construction du nom du vecteur est le même : \\

\begin{center}
  \textbf{<famille>\_<paramètre>\_<capteur>\_<type>\_<fréquence>.vbn}
\end{center}

avec :

\begin{center}
\begin{tabular}{|p{3cm}|p{11cm}|}
\hline

\textbf{<famille>} &
abréviation de la famille des données.

ex : \texttt{tpr} pour température.

Voir le fichier \texttt{\$AIDA\_HOME/config/familles}.
\\ \hline

\textbf{<paramètre>} &
mnémonique du nom physique de la grandeur mesurée.

ex : \texttt{air} pour une pression statique.
\\ \hline

\textbf{<capteur>} &
mnémonique du nom du (ou des) capteur(s).

ex : \texttt{li7500} pour un licor-7500.
\\ \hline

\textbf{<type>} &
type de données ou de taitement.

\begin{itemize}
  \item[\texttt{raw}] pour des données brutes.
  \item[\texttt{cal}] pour des données calibrées.
  \item[\texttt{xxx}] mnémonique de l’algorithme de calcul. 
                      Ex : \texttt{\textbf{Chb}} si une \textbf{C}orrection
                      des valeurs \textbf{h}ors\-\textbf{b}ornes a été appliquée.
\end{itemize}
\\ \hline

\textbf{<fréquence>} &
fréquence d’échantillonnage des donnée.

Les règles de formatage sont * :
\begin{itemize}
  \item[•] lorsque l’échantillonnage est irrégulier :

           la fréquence est écrite sous la forme 0.

           ex : \texttt{OMEGA\_\textbf{0}.vbn}

  \item[•] pour une fréquence supérieure à 2 Hz :

           la fréquence est écrite comme un nombre entier (arrondi au plus
           proche).

           ex : \texttt{ven\_v\_gillv\_cal\_\textbf{21}.vbn}

  \item[•] pour une fréquence ente 0.5 et 2 Hz :

           la fréquence est écrite comme un nombre décimal avec au plus
           2 chiffres après le point décimal.

           ex : \texttt{ctl\_tboacr2\_rs\_cal\_\textbf{0.83}.vbn}

  \item[•] pour une fréquence inférieure à 0.5 Hz (période supérieure à
           2 secondes) :

           on écrit la période (inverse de la fréquence), arrondie à l’entier
           le plus proche, précédée de «\%»

           ex : \texttt{pre\_air\_c0\_cal\_\textbf{\%60}.vbn}
\end{itemize}
\\ \hline

\end{tabular}
\end{center}

* Les règles de formatage font que la partie “fréquence” du nom de paramètre ne
sera pas systématiquement identique au champ fréquence de l’en-tête. Mais en
appliquant ces règles de formatage à la valeur contenue dans l’en-tête, on doit
obtenir une chaîne de caractère identique à la partie finale du nom.

